%% kapa: exam=1 round=36 year=2025 period=1 form=A subject=civil_law subject_ko=민법 seq=24 number=24
%% source: Q-Net 감정평가사 기출문제 (cst00309) artlSeq=5231525 file=raw/1008/5231525/2025년 제36회 감정평가사 1차시험 1교시 문제지 원본.pdf
%% conversion: kordoc 4.24.51 -> tools/gen1.py (auto); review=proofread-sonnet; flags=none
\begin{question}{24}
등기의 추정력에 관한 설명으로 옳지 않은 것은? (다툼이 있으면 판례에 따름)
\begin{choices}
  \item 소유권이전등기가 경료된 경우, 그 등기명의자는 전소유자에 대하여도 적법한 등기 원인에 의하여 소유권을 취득한 것으로 추정된다.
  \item 건물 소유권보존등기의 명의자가 그 건물을 신축한 것이 아니라면 그 등기의 권리 추정력은 깨진다.
  \item 소유권이전등기가 부적법하게 말소된 경우, 말소된 등기의 명의자는 여전히 적법한 소유자로 추정된다.
  \item 소유권이전등기청구권의 보전을 위한 가등기가 있는 경우, 소유권이전등기를 청구할 적법한 법률관계가 있다고 추정된다.
  \item 근저당권설정등기가 있는 경우에도 피담보채권의 성립을 위한 법률행위의 존재는 추정되지 않는다.
\end{choices}
\end{question}
\answer{4}
